\section{Split-Hadamard Information Maps}
\label{sec:split-hadamard}

We now identify a smooth global condition that produces rigid coherent
information sheets.  Let \((M,g)\) and \((B,h)\) be finite-dimensional
Hadamard manifolds, and let \(q:M\to B\) be a surjective Riemannian
submersion.  Its vertical and horizontal distributions are
\[
  \mathcal V_x=\ker Dq(x),
  \qquad
  \mathcal H_x=\mathcal V_x^{\perp_g}.
\]

\begin{definition}[Split-Hadamard information map]
\label{def:split-hadamard}
\mbox{}\par\noindent
The map \(q\) is split-Hadamard if \(\mathcal H\) is integrable and every
maximal horizontal leaf is complete and totally geodesic.
\end{definition}

This condition rules out holonomy between visible and invisible directions.
It does not require the fibers themselves to be linear or compact.

\begin{theorem}[Global split-Hadamard information theorem]
\label{thm:split-hadamard}
Let \(q:M\to B\) be split-Hadamard.  For every \(x\in M\), let \(L_x\) be
the horizontal leaf through \(x\).  Then:
\begin{enumerate}[label=(\roman*),leftmargin=*]
  \item The restriction \(q|_{L_x}:L_x\to B\) is a global Riemannian
  isometry.  Its inverse
  \[
    s_x:B\to L_x\subset M
  \]
  is the unique horizontal isometric totally geodesic section satisfying
  \(s_x(q(x))=x\).
  \item The map \(q\) is a metric submetry and \(s_x\) is the unique shortest
  completion:
  \begin{equation}
    s_x(b)=\argminop_{q(z)=b}d_M(x,z),
    \qquad
    d_M(x,s_x(b))=d_B(q(x),b).
    \label{eq:split-shortest-lift}
  \end{equation}
  \item Universal radial reduction \eqref{eq:universal-radial-reduction}
  holds.  Every reduced minimizer lifts through \(s_x\); if \(\rho\) is
  strictly increasing on its finite domain, these are all finite-valued full
  minimizers.
  \item Moreau envelopes and proximal maps commute exactly:
  \begin{align}
    e_\tau^M(\phi\circ q)(x)&=e_\tau^B\phi(q(x)),
    \label{eq:split-moreau}\\
    \prox^M_{\tau(\phi\circ q)}(x)
    &=s_x\!\left(\prox^B_{\tau\phi}(q(x))\right)
    \label{eq:split-prox}
  \end{align}
  for proper lower-semicontinuous geodesically convex bounded-below
  \(\phi\).  The visible proximal point is unique by Hadamard convexity; the
  full proximal point is unique because equality requires both that visible
  point and its unique shortest lift.
  \item If \(\phi\in C^1(B)\), then
  \begin{equation}
    \operatorname{grad}_M(\phi\circ q)(z)
    =
    \bigl(Dq(z)|_{\mathcal H_z}\bigr)^{-1}
    \operatorname{grad}_B\phi(q(z)).
    \label{eq:split-gradient}
  \end{equation}
  Consequently, every solution \(y:I\to B\) of
  \(\dot y=-\operatorname{grad}_B\phi(y)\) with \(y(0)=q(x)\) lifts on the
  same existence interval to the solution
  \(z=s_x\circ y\) of
  \(\dot z=-\operatorname{grad}_M(\phi\circ q)(z)\) with \(z(0)=x\).
  If \(\operatorname{grad}_B\phi\) is locally Lipschitz, the corresponding
  maximal solutions are unique.
\end{enumerate}
\end{theorem}

The proof in \cref{app:split-proofs} has two ingredients.  Completeness
extends every horizontal lift of a visible geodesic, and Hadamard uniqueness
makes the restriction to a horizontal leaf globally injective.  Equality in
Riemannian-submersion length contraction forces a shortest ambient geodesic
to be horizontal, yielding uniqueness of the nearest completion.

\begin{corollary}[Rigid coherent quotient]
\label{cor:split-coherent}
The sections \(s_x\) satisfy \eqref{eq:coherent-sections}; hence every
split-Hadamard map is a rigid coherent metric quotient in the sense of
\cref{thm:coherent-reduction}.
\end{corollary}

The abstract theorem is useful only after its hypotheses have been verified.
The next section proves them globally for SPD information compression and, in
addition, resolves the section in closed form.
